% Section 2: Related work.
\section{Related Work}
\label{sec:related}

\paragraph{Model judges and correlated output errors.}
Language models are widely used to grade outputs, provide feedback, and
supervise other models \citep{zheng2023judge,irving2018debate,bowman2022measuring,
burns2023weak}. This use assumes that a judge remains informative when the model
being judged is wrong. Several studies challenge that assumption at the output
level. \citet{panickssery2024selfpreference} show that model judges favor their
own generations. \citet{goel2025alike} introduce CAPA, a chance-adjusted measure
of agreement on predictions, and find that similarity affects both model
evaluation and weak-to-strong supervision. Across more than $350$ models,
\citet{kim2025correlated} find substantial conditional agreement on wrong
answers, while \citet{ding2026agree} show that agreement does not reliably imply
correctness. These works establish shared output error. \citet{kim2025correlated}
also analyze vendor and architecture effects, finding greater output-error
agreement within companies and model families. We ask the corresponding
directional question for internal readouts: whether cross-family selection
reduces false agreement once marginal error rates and model size are accounted
for. We also test whether shared errors can be detected from joint activations,
which can reveal truth-related information absent from a model's generated
response.

\paragraph{Truth probes and internal readouts.}
Linear probes can recover directions associated with truth from a model's
activations. \citet{burns2023ccs} propose an unsupervised consistency-based
method, and \citet{marks2024geometry} identify a linear truth direction with
causal relevance. Probes can also outperform a model's expressed confidence on
items where its output is wrong \citep{liu2023dissonance}, making them a
plausible instrument for oversight. Their interpretation nevertheless requires
care. Unsupervised methods may recover a salient binary feature other than truth
\citep{farquhar2024challenges}, and truth directions do not transfer uniformly
across tasks and logical forms \citep{bao2025probing}. We therefore validate the
probes against output confidence and base our main results on supervised linear
probes. Relatedly, \citet{pandey2026sycophancy} identify internal signals of a
false user claim across multiple models. Our unit of analysis is different. We
study the joint readouts of two models and the cases in which both fail in the same
direction.

\paragraph{Error correlation and ensemble diversity.}
That ensemble performance depends on covariance between member errors is
classical \citep{brown2005diversity}. Likewise, the Fr\'echet bounds describe
how the marginal rates of two binary events constrain their possible overlap
\citep{frechet1951tableaux}. We do not claim either result as new. We use them
to decompose the blind spot of probe-based oversight into the contribution of
individual error rates and the contribution of correlated error. This differs
from reporting a new similarity score. The decomposition identifies the range
of false agreement permitted by probe accuracy and explains why a reduction in
error correlation need not reduce the absolute overlap of errors.

\paragraph{Detecting disagreement and shared error.}
Cross-model disagreement can itself be useful. \citet{gorbett2026disagreement}
use one model's surprise at another's answer as a label-free correctness signal
and improve detection of confident errors over within-model uncertainty
baselines. Our disagreement router uses the analogous signal from two probe
verdicts. The harder case is common-mode error, where both readers agree and
there is no disagreement to route. \citet{afrin2026ensembles} show theoretically
that the common-mode component of evaluator error is not identifiable from the
evaluators' scores alone. We test an internal-state counterpart by asking whether false
agreement leaves a learnable signature in the models' activations, and whether
that signature transfers to an unseen data distribution.

\paragraph{What this paper adds.}
First, we define and measure false agreement between two models' internal
readouts rather than between their outputs. Second, we use an exact
decomposition and its bounds to distinguish the effects of probe accuracy from
error correlation, then measure both across model families, scales, datasets,
and probe estimators. Third, we test two proposed remedies: selecting a judge
from another model family and learning a detector for the residual blind spot.
The first does not reliably reduce false agreement, and the second reveals a
strong in-domain but largely domain-specific signature. Together, these results
locate the limit of cross-model agreement in shared error and make that limit a
measurable object for model-judge selection.
